\chapter*{Resumen}
El objetivo de hacer estos apuntes no es otro sino el de darles un texto de apoyo a las diapositivas difíciles de seguir de la asignatura, aunque pueden parecer un poco caóticos, ya que hay muchos saltos adelante y hacia atrás, si se siguen de manera paralela a las clases no se debe tener problema.\\\\
Los apuntes están divididos en Física de Partículas y Física Nuclear (la parte de fisión del tema 4 no está incluida pues no se impartió en el curso académico 24-25 al igual que el final del tema 5), además de una serie de Apéndices con teoría de Mecánica Cuántica que se escapa del currículo de la asignatura, pero que se suponen conocidas por parte del profesor, aún así, se abordan estos temas en clase. Estos apuntes no están libres de erratas, pues aunque los he revisado junto con otros compañeros es mucho contenido y se pueden escapar algunas pinceladas.\\\\

\begin{flushright}
    Mucha suerte a todos y espero que os sirvan por mucho tiempo.
\end{flushright}

\begin{flushright}
    Chinri
\end{flushright}

\vfill
\begin{flushleft}
        % Nota de derechos y autoría
    \noindent\textbf{Licencia y Derechos de Autor} \\
    \vspace{0.2cm}
    \small
    Este documento y su código fuente son de elaboración propia. 
    Se distribuyen bajo una licencia \\ \textbf{Creative Commons Attribution-NonCommercial-ShareAlike 4.0 International} (CC BY-NC-SA 4.0).
    
    \vspace{0.4cm}
    
    % Esto inserta los iconos típicos de Creative Commons de forma limpia usando texto
    \includegraphics[width=0.25\textwidth]{by-nc-sa.pdf}
    
    \vspace{0.4cm}
    \footnotesize
    Eres libre de compartir y adaptar este material siempre que se reconozca la autoría original, 
    no se utilice con fines comerciales y se distribuya bajo la misma licencia.
\end{flushleft}